\section{Ramification of Discrete Valuations}\label{section:ramification_of_discrete_valuations}
This section reviews the ramification of discrete valuations, beginning with the general theory of \textit{discrete valuation rings} and their
integral closures in finite separable extensions.
We then specialize to complete rings in Galois extensions to describe the ramification filtration via lower and upper numbering,
following the treatments in \stackstag{0EXQ} and \cite{serreLF}.

\begin{definition}[\stackstag{09E4}]\label{definition:weakly_unramified_and_tamely_ramified_DVR}
    We say that $A \to B$ or $A \subset B$ is an extension of discrete valuation rings if $A$ and $B$ are discrete valuation rings and $A \to B$ is injective and local. In particular, if $\pi _ A$ and $\pi _ B$ are uniformizers of $A$ and $B$, then $\pi _ A = u \pi _ B^ e$ for some $e \geq 1$ and unit $u$ of $B$. The integer $e$ does not depend on the choice of the uniformizers as it is also the unique integer $\geq 1$ such that
\[ \mathfrak m_ A B = \mathfrak m_ B^ e \]
The integer $e$ is called the \textit{ramification index} of $B$ over $A$. We say that $B$ is \textit{weakly unramified} over $A$ if $e = 1$.
If the extension of residue fields $\kappa _ A = A/\mathfrak m_ A \subset \kappa _ B = B/\mathfrak m_ B$ is finite,
then we set $f = [\kappa _ B : \kappa _ A]$ and we call it the \textit{residual degree} or residue degree of the extension $A \subset B$.

Note that we do not require the extension of fraction fields to be finite.
\end{definition}

Now let $A$ be a discrete valuation ring with fraction field $K$, let  $L/K$ be a finite separable field extension and
let $B \subset L$ be the integral closure of $A$ in $L$.
Then the ring extension $A \subset B$ is finite, hence $B$ is Noetherian.
The dimension of $B$ is $1$, hence $B$ is a Dedekind domain.
Let $\mathfrak m_1, \ldots , \mathfrak m_ n$ be the maximal ideals of $B$ (i.e., the primes lying over $\mathfrak m_ A$).
We obtain extensions of discrete valuation rings

\[ A \subset B_{\mathfrak m_ i} \]
and hence ramification indices $e_ i$ and residue degrees $f_ i$. We have

\[ [L : K] = \sum \nolimits _{i = 1, \ldots , n} e_ i f_ i \]
Note that if $A$ is henselian (e.g. $A$ is complete) then $n=1$.

\begin{definition}[\stackstag{09E9}]\label{definition:tamely_ramified_dvrs}
Let $A$ be a discrete valuation ring with fraction field $K$. Let $L/K$ be a finite separable extension. With $B$ and $\mathfrak m_ i$, $i = 1, \ldots , n$
as above, we say the extension $L/K$ is
\begin{enumerate}
    \item unramified with respect to $A$ if $e_ i = 1$ and the extension $\kappa (\mathfrak m_ i)/\kappa _ A$ is separable for all $i$,
    \item tamely ramified with respect to $A$ if, for every $i$, the residue field extension
    $\kappa(\mathfrak m_i)/\kappa_A$ is separable and the ramification index $e_i$ is prime to
    the residue characteristic. Thus, when $\operatorname{char}(\kappa_A)=0$, the condition on
    $e_i$ is vacuous (and the residue field extensions are automatically separable), whereas when
    $\operatorname{char}(\kappa_A)=p>0$, each $e_i$ must be prime to $p$, and
    \item totally ramified with respect to $A$ if $n = 1$ and the residue field extension $\kappa (\mathfrak m_1)/\kappa _ A$ is trivial.
\end{enumerate}
If the discrete valuation ring $A$ is clear from context, then we sometimes say $L/K$ is unramified, totally ramified, or tamely ramified for short.
\end{definition}

For the remainder of this section, assume that $A$ is a discrete valuation ring of equal characteristic
$p>0$, complete with respect to its maximal-ideal-adic topology (cf.\ \stackstag{00IE} and
\stackstag{0324}), with uniformizer $\pi$ and perfect residue field $\kappa$. Thus,
$\operatorname{char}(A)=\operatorname{char}(\kappa)=p$. The Cohen structure theorem
(\stackstag{032A}; see also \stackstag{0C0S}) provides a coefficient field $k\subset A$ mapping
isomorphically onto $\kappa$. Consequently, $A\cong k[[\pi]]\cong k[[t]]$, and its fraction field is
$K\cong k((\pi))$.
Let $L/K$ be a finite Galois extension with Galois group $G$.
The integral closure $B$ of $A$ in $L$ is itself a complete discrete valuation ring.
Since the residue field $\kappa$ is perfect, the extension of residue fields is separable, and there exists a uniformizer
$\Pi \in B$ such that $B = A[\Pi]$.

The \textit{lower numbering ramification filtration} of $G$, denoted by $(G_i)_{i \geq -1}$, is defined as:
\[
G_i = \{ \sigma \in G \mid v_B(\sigma (x) - x) \geq i + 1 \text{ for all } x \in B \}
\]
where $v_B$ is the valuation on $L$ associated with $B$. In this indexing, $G_{-1} = G$ and $G_0$
corresponds to the \textit{inertia group} of the extension $L/K$. Note that $L/K$ is unramified if and only if
$G_0 = \{1\}$, and it is tamely ramified if and only if $G_1 = \{1\}$.

A standard result in the theory of local fields (complete local fields with a discrete valuation and perfect residue field)
ensures that the condition in the definition of $G_i$ need only be checked for the uniformizer $\Pi$ of $B$.
Specifically, if we define the function
\[
i^L_K(\sigma) = v_B(\sigma(\Pi) - \Pi)
\]
for $\sigma \in G$, then $G_i = \{ \sigma \in G \mid i^L_K(\sigma) \geq i + 1 \}$.
The subgroups $G_i$ are normal in $G$ and become trivial for sufficiently large $i$.

Consider a tower of fields $K \subset E \subset L$, and let $H = \text{Gal}(L/E) \subset G$. The lower numbering is compatible with subgroups:
\[
G_i \cap H = H_i \quad \text{for all } i \geq -1
\]
which follows from the identity $i^L_E = i^L_K|_{H}$.
However, the lower numbering does not behave as simply with respect to quotients. If $H$ is normal in $G$, the image of $G_i$ in $G/H$ is generally not $(G/H)_i$.
Instead, we have $G_i H / H = (G/H)_j$, where the index $j$ is given by the formula:
\[
j = \frac{1}{e_{L/E}} \sum_{\tau \in H} \min(i^L_K(\tau), i+1) - 1
\]


In the literature, one reindexes the ramification groups by defining the Herbrand function $\phi_{L/K} : [-1, \infty ) \to [-1, \infty )$:
\[ \phi_{L/K}(i) = \frac{1}{e_{L/K}}\sum_{\sigma \in G} \min(i^L_K(\sigma), i+1)-1 = \int _0^i \frac{1}{[G_0 : G_t]} dt \]

This function is continuous, strictly increasing, and piecewise linear, making it a bijection.
It satisfies the composition law $\phi_{L/K} = \phi_{E/K} \circ \phi_{L/E}$ for a tower of extensions $K \subset E \subset L$.
Using this function, we define the \textit{upper numbering ramification groups} as:
\[
G^i = G_{\phi_{L/K}^{-1}(i)}
\]
The primary advantage of this reindexing is its compatibility with quotients: for any normal subgroup $H \subset G$, we have:
\[
G^i H / H = (G/H)^i
\]
for all $i \geq -1$.

We record, for orientation, the structure of the inertia group and of the
graded pieces of the filtration; these facts are not used in the sequel, but
they explain the dichotomy --- prime-to-$p$ versus $p$-power --- that shapes
the whole chapter.

\begin{prop}[\Cite{serreLF}, IV Corollary 3]
    The quotients $G_i/G_{i+1}$, $i \geq 1$, are abelian groups, and are direct products of cyclic groups of order $p$.
    The group $G_1$ is a $p$-group.
\end{prop}

\begin{prop}[\Cite{serreLF}, IV Corollary 4]
    The inertia group $G_0$ has the following property:
    \textit{It is the semi-direct product of a cyclic group of order prime to $p$ with a normal subgroup whose order is a power of $p$.}
\end{prop}
